\section{Conclusion}

We have presented an empirical characterization of an 11-stage AI-assisted pipeline for puzzle hunt creation, executed across five independent scenarios spanning game-knowledge, noir fiction, science-fiction RPG, music, and investigation domains. The pipeline produced 52 puzzles across these scenarios, of which 47 passed editorial review on the first pass — a 90.4\% first-pass pass rate — with all five failures recoverable by Stage~11. To our knowledge, this is the first cross-scenario empirical characterization of a structured generative pipeline for multi-artifact creative systems.

Three findings from this study are robust across the full scenario set. First, the 90.4\% first-pass pass rate, held consistently across five domains and two scales (5--19 puzzles), establishes that the pipeline produces shippable output as a reliable baseline rather than as an occasional outcome. Unconstrained generation of creative artifacts at this level of quality is not achievable without the stage structure the pipeline imposes; the structure is the intervention, not the individual prompts. Second, the concentration of failures at Stage~8 (testing and integration), rather than at the review gates (Stages~3 and~7), confirms that the gates are functioning as designed. A gate that catches everything would leave nothing for downstream stages to fail; a gate that catches nothing would redistribute failures uniformly across stages. The empirical distribution — zero gate-stage failures, three Stage~8 failures — is consistent with gates operating at the right threshold: permissive enough to let strong candidates through, selective enough to filter weak ones before authoring investment. Third, the same skill library — 13 \texttt{CLAUDE.md}-driven skills and a 29-profile reviewer panel — executed all five scenarios without per-scenario modification. Only the scope and world documents varied. This confirms cross-theme, cross-setting transferability at the specification level, which is the form of transferability that matters for practical adoption.

For practitioners considering pipeline adoption, the central message is that structure is the primary quality lever. Individual prompt quality, reviewer profile sophistication, and domain-specific knowledge all contribute, but the two review gates at Stages~3 and~7 account for the majority of the quality guarantee. The world lock — a single committed document that all subsequent stages treat as read-only — provides coherence insurance at negligible cost. These are structural decisions, not implementation details, and they are the first thing a practitioner should get right.

Several questions remain open. Pipeline behavior at MIT Mystery Hunt scale — 50 or more puzzles with nested meta structure and cross-round answer dependencies — is the most pressing empirical gap; the coherence properties of the world lock and the compounding probability of Stage~8 coordination failures both warrant investigation at that scale. The correlation between AI panel editorial verdict and human solver experience is studied in a companion paper~\cite{dellaLibera2025calibration}; integrating those results with the stage-level pass rates reported here is a natural next step. Cross-setting pipeline adaptation — to tabletop campaign generation, interactive fiction, or educational puzzle design — would test whether the specification-level transferability observed across five puzzle hunt scenarios spanning different themes and settings extends to structurally different creative artifact classes. Finally, the failure data collected across these five scenarios provides a natural training signal for automated skill refinement: failure types can be mapped to the stage and rubric criterion that should have caught them earlier, enabling systematic improvement of the review gates over time.

The full pipeline specification, including the \texttt{CLAUDE.md} skill library, reviewer profiles, and scenario artifacts for all five scenarios, is available at [URL]. We offer it as a reusable foundation for researchers and practitioners working on structured generative pipelines for complex creative systems.
